% =====================================================================
%  Webseite Förderverein KF/JF Kreuzau -- Anleitung für Mitglieder
%  Vorlage: deb-support (debsupport.sty, siehe latex/README.md)
%  Bilder: Bildschirmfotos einer Testinstallation mit Beispieldaten
%  Übersetzen mit:  pdflatex anleitung_mitglieder_webseite.tex  (2x)
% =====================================================================
\documentclass[a4paper,fontsize=10pt]{scrartcl}
\usepackage[bericht]{debsupport}
\usepackage{float}                % Bilder genau an ihrer Stelle
\setcounter{tocdepth}{1}
\renewcommand*{\textfraction}{0.1}
\renewcommand*{\topfraction}{0.9}
\renewcommand*{\floatpagefraction}{0.8}

\renewcommand*{\dsdokumentart}{Anleitung}
\renewcommand*{\dskurztitel}{Webseite Förderverein KF/JF Kreuzau -- Anleitung für Mitglieder}
\renewcommand*{\dsbeleg}{}
\setlength{\emergencystretch}{3em}
\graphicspath{{bilder/}}

% Bildschirmfoto mit dünnem Rahmen:  \foto[Breite]{Datei}{Bildunterschrift}
\newcommand{\foto}[3][0.72\textwidth]{%
  \begin{figure}[H]
    \centering
    \fcolorbox{dsgrau!60}{white}{\includegraphics[width=#1]{#2}}
    \caption{#3}
  \end{figure}}
% Bedienelement, z. B. \knopf{Anmelden}
\newcommand*{\knopf}[1]{\textbf{\guillemotright #1\guillemotleft}}
\newcommand*{\web}{\url{https://foerderverein-kf-jf-kreuzau.de}}
\setcapindent{0pt}\addtokomafont{caption}{\small}\addtokomafont{captionlabel}{\bfseries}

\begin{document}

\dstitel{Die Webseite des Fördervereins}
        {Anleitung für Mitglieder und Fördermitglieder:
         Konto, Anmelden, Seiten und Anhänge, Newsletter}
        {Webseite: & \web\\
         Für:      & alle Personen mit Konto (ohne Bearbeitungsrechte)\\
         Stand:    & \today}

\begin{dshinweis}[title=Kurz und knapp]
\begin{enumerate}
  \item Sie erhalten vom Vorstand eine \textbf{Einladung per E-Mail}. Über den
        Link darin legen Sie Ihr Passwort fest.
  \item Danach melden Sie sich mit \textbf{Benutzername oder E-Mail-Adresse}
        und Passwort an (\knopf{Anmelden} oben rechts).
  \item Auf der \textbf{Startseite} stehen alle Beiträge, die Sie sehen dürfen,
        die neuesten zuerst. \knopf{Weiterlesen} öffnet einen Beitrag mit
        Bildern und Anhängen.
  \item Unter \textbf{Ihrem Namen} oben rechts finden Sie \emph{Mein Konto}:
        Passwort ändern und weitere Einstellungen.
\end{enumerate}
\end{dshinweis}

\tableofcontents

% ------------------------------------------------------------------
\section{Was bietet die Webseite?}

Die Webseite hat einen öffentlichen Teil, den jeder ohne Anmeldung sieht,
und einen geschützten Teil für Mitglieder. Nach der Anmeldung sehen Sie
zusätzlich die Beiträge für Ihre Gruppe: Einladungen, Protokolle,
Informationen und die zugehörigen Dateien (meist PDF).

Jedes Konto gehört zu einer \textbf{Stufe}. Wer zu einer Stufe gehört, sieht
deren Beiträge und die aller niedrigeren Stufen:

\begin{center}
\begin{tabularx}{0.85\textwidth}{@{}lX@{}}
\toprule
\textbf{Stufe} & \textbf{sieht} \\
\midrule
Öffentlich (ohne Anmeldung) & öffentliche Beiträge \\
Fördermitglieder            & dazu Beiträge für Fördermitglieder \\
Mitglieder                  & dazu Beiträge für Mitglieder \\
Vorstand                    & alles; der Vorstand schreibt auch die Beiträge \\
\bottomrule
\end{tabularx}
\end{center}

Geschützte Beiträge erkennen Sie an einem kleinen Kästchen mit dem Namen der
Stufe neben dem Datum (z.\,B.\ \fbox{\small Mitglieder}).
Öffentliche Beiträge haben kein solches Kästchen.

Die Webseite funktioniert in jedem aktuellen Browser auf PC, Tablet und
Smartphone. Eine App ist nicht nötig.

% ------------------------------------------------------------------
\section{Ihr Konto einrichten}

\subsection{Die Einladung}

Konten legt der Vorstand an; selbst registrieren kann man sich nicht. Sobald
Ihr Konto angelegt ist, erhalten Sie eine E-Mail mit dem Betreff
\emph{[Förderverein KF/JF Kreuzau] Einladung zur Webseite}:

\begin{dshinweis}[title=Beispiel: Einladung]
\small
Hallo Erika Neu,\par
für Sie wurde ein Konto auf der Webseite \glqq Förderverein KF/JF Kreuzau\grqq{}
angelegt: https://foerderverein-kf-jf-kreuzau.de\par
Ihr Benutzername: erika.neu\\
(Anmelden können Sie sich auch mit dieser E-Mail-Adresse.)\par
Bitte legen Sie über den folgenden Link Ihr Passwort fest. Der Link ist
7 Tage gültig und kann nur einmal verwendet werden:\\
https://foerderverein-kf-jf-kreuzau.de/konto/password/reset/key/…\par
Mit freundlichen Grüßen\\
Der Vorstand
\end{dshinweis}

Notieren Sie sich Ihren \textbf{Benutzernamen}. Sie können sich aber auch
immer mit Ihrer E-Mail-Adresse anmelden.

\subsection{Passwort festlegen}

Ein Klick auf den Link in der Einladung öffnet die Seite \emph{Passwort
ändern}. Geben Sie dort zweimal dasselbe neue Passwort ein und klicken Sie
auf \knopf{Passwort ändern}.

\foto{02_passwort_festlegen}{Passwort festlegen über den Link aus der Einladung}

Das Passwort muss mindestens 10~Zeichen lang sein, darf nicht nur aus Ziffern
bestehen und nicht zu einfach sein (z.\,B.\ nicht \emph{passwort123} oder Ihr
Name). Gut merken lässt sich ein Satz oder mehrere Wörter, etwa
\emph{Sommerfest-am-Gerätehaus-2027}.

\begin{dshinweis}[title=Link abgelaufen?]
Der Link gilt 7~Tage und nur einmal. Ist er abgelaufen, fordern Sie über
\emph{Passwort vergessen?} auf der Anmeldeseite einfach einen neuen an
(Abschnitt~\ref{sec:vergessen}).
\end{dshinweis}

% ------------------------------------------------------------------
\section{Anmelden und Abmelden}

Klicken Sie auf der Webseite oben rechts auf \knopf{Anmelden}.

\foto{03_anmelden}{Die Anmeldeseite}

\begin{enumerate}
  \item Bei \emph{Anmeldung} tragen Sie Ihren \textbf{Benutzernamen} (z.\,B.\
        \emph{max.mustermann}) oder Ihre \textbf{E-Mail-Adresse} ein.
  \item Bei \emph{Passwort} tragen Sie Ihr Passwort ein.
  \item \emph{Angemeldet bleiben} sollten Sie nur an Ihrem eigenen Gerät
        ankreuzen. Dann müssen Sie sich dort nicht jedes Mal neu anmelden.
  \item Klicken Sie auf \knopf{Anmelden}.
\end{enumerate}

Nach der Anmeldung steht oben rechts Ihr Name. Ein Klick darauf öffnet
\emph{Mein Konto}. Daneben steht \knopf{Abmelden}. Melden Sie sich an fremden
oder gemeinsam genutzten Geräten nach dem Lesen immer ab.

\foto{04b_kopfzeile}{Kopfzeile nach der Anmeldung: Ihr Name (Mein Konto) und Abmelden}

% ------------------------------------------------------------------
\section{Die Startseite}

Die Startseite zeigt alle Beiträge, die Sie sehen dürfen, die neuesten oben.
Jeder Beitrag erscheint mit Datum, Titel, den ersten Zeilen des Textes und
gegebenenfalls einem Vorschaubild.

\foto{04_startseite_mitglied}{Startseite nach der Anmeldung als Mitglied}

\begin{itemize}
  \item Das Kästchen neben dem Datum (hier \fbox{\small Mitglieder}) zeigt,
        für wen der Beitrag bestimmt ist.
  \item \emph{2~Anhänge} bedeutet: Zum Beitrag gehören zwei Dateien, z.\,B.\
        eine Einladung als PDF.
  \item \knopf{Weiterlesen} oder ein Klick auf den Titel öffnet den ganzen
        Beitrag.
  \item Ein Klick auf den Vereinsnamen oben links bringt Sie von überall
        zurück zur Startseite.
\end{itemize}

Ohne Anmeldung sieht die Startseite genauso aus, enthält aber nur die
öffentlichen Beiträge:

\foto[0.6\textwidth]{01_startseite_gast}{Startseite ohne Anmeldung: nur öffentliche Beiträge}

% ------------------------------------------------------------------
\section{Einen Beitrag lesen}

\subsection{Text und Anhänge}

Ein geöffneter Beitrag zeigt den vollständigen Text. Darunter stehen unter
\emph{Anhänge} die zugehörigen Dateien. Ein Klick auf einen Anhang öffnet
oder speichert die Datei, je nach Browser.

\foto[0.65\textwidth]{05_seite_jhv}{Ein Beitrag für Mitglieder mit zwei Anhängen}

\begin{dshinweis}[title=Anhänge sind geschützt]
Anhänge eines geschützten Beitrags lassen sich nur abrufen, wenn Sie
angemeldet sind und den Beitrag sehen dürfen. Ein weitergeleiteter Link auf
einen Anhang funktioniert für andere Personen daher nicht. Bitte geben Sie
heruntergeladene Dateien aus dem Mitgliederbereich nicht an Außenstehende
weiter.
\end{dshinweis}

\subsection{Bilder}

Hat ein Beitrag Bilder, steht das \textbf{Titelbild} groß über dem Text, die
übrigen Bilder stehen als kleine \textbf{Galerie} darunter. Ein Klick auf
ein Bild zeigt es in voller Größe; mit der Zurück-Taste des Browsers kommen
Sie zum Beitrag zurück.

\foto[0.55\textwidth]{06_seite_fest}{Beitrag mit Titelbild, Text und Galerie
  (Beispielbilder)}

% ------------------------------------------------------------------
\section{Mein Konto}

Ein Klick auf Ihren Namen oben rechts öffnet \emph{Mein Konto}. Dort sehen
Sie Ihren Namen, Benutzernamen, Ihre E-Mail-Adresse und Ihre Stufe.

\foto{07_mein_konto}{Mein Konto}

Name, E-Mail-Adressen und Stufe ändert der Vorstand. Hat sich zum Beispiel
Ihre E-Mail-Adresse geändert, sagen Sie bitte dem Vorstand Bescheid. Sonst
erreichen Sie Einladungen und \emph{Passwort vergessen} nicht mehr.

\subsection{Vereinsadresse und Vereinsmails}

Einige Personen haben zusätzlich eine \textbf{Vereinsadresse} (ein Postfach
\emph{…@foerderverein-kf-jf-kreuzau.de}). Dann erscheint auf \emph{Mein Konto}
der Abschnitt \emph{Vereinsmails}. Dort wählen Sie, wohin E-Mails des
Vereins, etwa der Newsletter, geschickt werden: an Ihre externe Adresse oder
an die Vereinsadresse. Auswahl treffen und \knopf{Speichern} klicken.

\foto[0.65\textwidth]{07b_konto_vereinsmails}{Mein Konto mit Vereinsadresse:
  Auswahl, wohin Vereinsmails gehen}

Mit der Vereinsadresse können Sie sich auch an der Webseite anmelden. Das
Passwort ist für Webseite und Postfach dasselbe.

% ------------------------------------------------------------------
\section{Passwort}

\subsection{Passwort ändern}

Unter \emph{Mein Konto} $\rightarrow$ \knopf{Passwort ändern} geben Sie Ihr
bisheriges Passwort und zweimal das neue ein.

\foto[0.7\textwidth]{08_passwort_aendern}{Passwort ändern}

Haben Sie eine Vereinsadresse, gilt das neue Passwort auch für Ihr Postfach.
Tragen Sie es danach in Ihren Mailprogrammen (Smartphone, PC) neu ein.

\subsection{Passwort vergessen}\label{sec:vergessen}

\begin{enumerate}
  \item Auf der Anmeldeseite auf \knopf{Passwort vergessen?} klicken.
  \item Ihre E-Mail-Adresse eintragen und \knopf{Passwort zurücksetzen}
        klicken.
  \item Sie erhalten eine E-Mail mit einem Link. Über diesen Link legen Sie
        wie bei der Einladung ein neues Passwort fest.
\end{enumerate}

\foto[0.7\textwidth]{09_passwort_vergessen}{Passwort vergessen: E-Mail-Adresse eintragen}

Kommt keine E-Mail an, sehen Sie bitte im Spam-Ordner nach. Aus Sicherheitsgründen
zeigt die Webseite immer dieselbe Bestätigung, auch wenn die Adresse nicht bekannt
ist. Hilft das nicht, wenden Sie sich an den Vorstand.

% ------------------------------------------------------------------
\section{Zwei-Faktor-Anmeldung (freiwillig)}

Für Mitglieder ist die Zwei-Faktor-Anmeldung freiwillig, für den Vorstand
Pflicht. Sie schützt Ihr Konto zusätzlich: Bei der Anmeldung braucht man
außer dem Passwort einen sechsstelligen Code aus einer App auf Ihrem
Smartphone. Ein gestohlenes Passwort allein reicht dann nicht mehr.

Sie benötigen eine Authenticator-App, zum Beispiel \emph{Aegis} oder
\emph{FreeOTP} (frei), \emph{Google Authenticator} oder \emph{Microsoft
Authenticator}.

\begin{enumerate}
  \item \emph{Mein Konto} $\rightarrow$ \emph{Zwei-Faktor-Anmeldung}
        $\rightarrow$ \knopf{einrichten}.
  \item In der App ein neues Konto hinzufügen und den QR-Code scannen.
  \item Den sechsstelligen Code aus der App bei \emph{Authentifizierungscode}
        eintragen und \knopf{Aktivieren} klicken.
\end{enumerate}

\foto[0.6\textwidth]{10_zwei_faktor}{Zwei-Faktor-Anmeldung einrichten}

Ab jetzt fragt die Webseite nach dem Passwort zusätzlich nach dem aktuellen
Code aus der App.

\begin{dshinweis}[title=Wichtig: für den Notfall vorsorgen]
Unter \emph{Mein Konto} $\rightarrow$ \emph{Zwei-Faktor-Anmeldung}
$\rightarrow$ \knopf{verwalten} können Sie \textbf{Wiederherstellungscodes}
erzeugen. Drucken Sie diese aus oder notieren Sie sie und bewahren Sie sie
sicher auf. Mit einem dieser Codes kommen Sie auch ohne Smartphone in Ihr
Konto, etwa bei Verlust oder Wechsel des Geräts.
\end{dshinweis}

% ------------------------------------------------------------------
\section{Newsletter per E-Mail}

Der Vorstand kann Beiträge zusätzlich als E-Mail an alle Personen schicken,
die den Beitrag sehen dürfen, zum Beispiel die Einladung zur
Jahreshauptversammlung. Die E-Mail kommt an die Adresse, die Sie unter
\emph{Vereinsmails} gewählt haben, sonst an Ihre externe Adresse.

Eine solche E-Mail enthält oft eine kurze Ankündigung, darunter entweder den
ganzen Beitrag oder nur den Anfang mit dem Knopf \knopf{Weiterlesen}.
Beiträge für Mitglieder öffnen sich erst nach der Anmeldung; danach landen
Sie direkt beim Beitrag.

\foto[0.6\textwidth]{11_newsletter_mail}{So sieht ein Newsletter im Postfach aus}

Wenn Sie auf den Newsletter antworten, geht Ihre Antwort an die Person, die
ihn verschickt hat, und an den Vorstand.

% ------------------------------------------------------------------
\section{Auf dem Smartphone}

Auf dem Smartphone öffnen Sie die Webseite einfach im Browser. Die Seite
passt sich der Bildschirmbreite an; Anmelden, Lesen und Anhänge funktionieren
wie am PC.

\foto[0.32\textwidth]{12_handy_startseite}{Startseite auf dem Smartphone}

Tipp: Über das Menü des Browsers (\emph{Zum Startbildschirm hinzufügen})
legen Sie die Webseite wie eine App auf den Startbildschirm.

% ------------------------------------------------------------------
\section{Häufige Fragen}

\begin{description}[style=nextline,leftmargin=0pt,itemsep=0.6\baselineskip]
  \item[Ich habe meinen Benutzernamen vergessen.]
    Melden Sie sich einfach mit Ihrer E-Mail-Adresse an.
  \item[Der Link aus der Einladung funktioniert nicht.]
    Er ist nach 7~Tagen oder nach einmaliger Verwendung ungültig. Fordern Sie
    über \emph{Passwort vergessen?} einen neuen an.
  \item[Ein Link aus einer E-Mail zeigt „nicht gefunden“.]
    Der Beitrag ist für eine andere Stufe bestimmt oder wurde entfernt. Prüfen
    Sie, ob Sie angemeldet sind. Sonst fragen Sie den Vorstand.
  \item[Ich sehe einen Beitrag nicht, den andere sehen.]
    Er ist wahrscheinlich für eine höhere Stufe bestimmt. Ob Ihre Stufe
    stimmt, sehen Sie unter \emph{Mein Konto}.
  \item[Ich bekomme keine Newsletter.]
    Prüfen Sie den Spam-Ordner und unter \emph{Mein Konto}, an welche Adresse
    Vereinsmails gehen. Stimmt die Adresse nicht, sagen Sie dem Vorstand
    Bescheid.
  \item[Kann ich den Newsletter abbestellen?]
    Nein. Über den Newsletter verschickt der Verein Mitteilungen an seine
    Mitglieder, zum Beispiel Einladungen zur Mitgliederversammlung.
  \item[Smartphone mit der Authenticator-App verloren.]
    Melden Sie sich mit einem Wiederherstellungscode an und richten Sie die
    App neu ein. Haben Sie keine Codes, wenden Sie sich an den Vorstand.
\end{description}

% ------------------------------------------------------------------
\section{Hilfe}

Bei Fragen zu Ihrem Konto, zu Ihrer Stufe oder zu Ihren E-Mail-Adressen
wenden Sie sich bitte an den Vorstand des Fördervereins.

{\small\color{dsgrau} Die Bildschirmfotos stammen aus einer Testinstallation
mit Beispieldaten; Personen, Beiträge und Bilder sind erfunden.}

\end{document}
